\chapter{Hoy: lo que necesita tu atención}\label{cap:hoy}

\begin{objbox}
Entenderás la pantalla de inicio: qué resume, qué puedes hacer desde ella y cómo te ayuda a empezar el día sin abrir diez ventanas.
\end{objbox}

\section{Para qué sirve}
\ui{Hoy} responde a una pregunta: \emph{¿qué tengo que hacer ahora?} Reúne en una sola pantalla los mensajes por contestar, los cobros pendientes, las tareas del día y las propuestas del asistente que esperan tu aprobación.

\figura[0.66]{e01-hoy}{Pantalla Hoy con datos ficticios de Estudio Marea Alta.}
\begin{leyenda}
\item[\mk{1}] \textbf{Buscar} (\cmd\,K). Abre la paleta de búsqueda y órdenes.
\item[\mk{2}] \textbf{Nuevo}. Menú de creación rápida: mensaje pegado, factura, presupuesto, cliente, gasto, tarea.
\item[\mk{3}] \textbf{Barra del asistente}. Se usa de dos formas: una \emph{pregunta} («¿cuánto cuesta el desplazamiento?»), que se contesta con los documentos de tu negocio, o una \emph{orden} («factura a Marta por 3 horas a 40 €»), que primero te enseña lo que va a hacer y solo lo hace si confirmas.
\item[\mk{4}] \textbf{Revisar mensajes}. Lleva a la bandeja con los mensajes por responder.
\item[\mk{5}] \textbf{Reclamar cobros}. Lleva a las facturas vencidas para reclamarlas.
\item[\mk{6}] \textbf{Tarjetas de cifras}: \emph{Por cobrar} (suma de facturas abiertas), \emph{Vencido} (la parte que ya pasó de fecha) y \emph{Este mes} (ingresos menos gastos del mes).
\item[\mk{7}] \textbf{Reclamar}. En «Necesita tu atención», prepara un recordatorio de pago para una factura vencida. No se envía: se abre un borrador que tú revisas.
\item[\mk{8}] \textbf{Facturar}. Aparece junto a un presupuesto aceptado y lo convierte en un borrador de factura.
\item[\mk{9}] \textbf{Casilla de tarea}. Marca una tarea como hecha sin salir de Hoy.
\item[\mk{10}] \textbf{Mensajes} en la barra lateral, con la cuenta de los pendientes.
\end{leyenda}

\section{Los bloques de la pantalla}
\begin{description}[leftmargin=0.6cm,style=nextline]
\item[Cabecera.] Un saludo con el resumen en una frase: facturas vencidas, tareas para hoy, presupuestos aceptados sin facturar y mensajes por responder.
\item[Mensajes por responder.] Los últimos mensajes de tus clientes, con el canal de origen. Al pulsar uno entras en su conversación.
\item[Tareas.] Las tareas de hoy y las atrasadas, ordenadas por prioridad y fecha.
\item[Necesita tu atención.] Facturas vencidas, presupuestos aceptados por facturar y propuestas del asistente. Las propuestas tienen dos botones: \ui{Aprobar} y \ui{Descartar}.
\item[Primeros pasos.] Una lista de cinco o seis tareas de puesta en marcha (datos del negocio, primer cliente, primera factura, copia de seguridad…). Se oculta con \ui{Ocultar primeros pasos}.
\end{description}

\begin{ideabox}
En REM, \emph{aprobar no es enviar}. El asistente puede dejar preparados un recordatorio o una respuesta, pero sale de tu equipo solo cuando tú pulsas el botón de enviar. Esa separación se explica en el capítulo~\ref{cap:asistente}.
\end{ideabox}

\begin{checkbox}
Escribe en la barra del asistente \texttt{tarea llamar al proveedor mañana}. Verás una ventana que dice qué tarea va a crear y para qué fecha. Cancela o confirma: hasta que no confirmas, no existe.
\end{checkbox}
